% =============================================================================
% Chapter: Posterior Inference for Treatment Effects
%
% Drop-in chapter for the data-fusion working document. Self-contained
% (report class + manual bibliography) so it can be compiled in isolation,
% but the chapter body can also be \input into the main document, in which
% case the preamble and \begin{document}/\end{document} should be removed
% and the \bibitem entries merged into the master bibliography.
%
% Primary estimands: the conditional average treatment effect (CATE) and the
% acceleration factor AF = exp(CATE).
%
% Not compiled.
% =============================================================================

\documentclass[11pt]{report}

\usepackage[margin=1.1in]{geometry}
\usepackage{amsmath,amssymb,amsthm}
\usepackage{bm}
\usepackage{booktabs}
\usepackage{natbib}
\usepackage{enumitem}

\theoremstyle{plain}
\newtheorem{proposition}{Proposition}[chapter]
\theoremstyle{definition}
\newtheorem{definition}{Definition}[chapter]
\newtheorem{remark}{Remark}[chapter]

% --- notation shortcuts -------------------------------------------------------
\newcommand{\E}{\mathbb{E}}
\newcommand{\Var}{\mathrm{Var}}
\newcommand{\Prob}{\mathbb{P}}
\newcommand{\given}{\,\vert\,}
\newcommand{\data}{\mathcal{D}}
\newcommand{\X}{\mathcal{X}}
\newcommand{\af}{\mathrm{AF}}
\newcommand{\vte}{\mathrm{VTE}}
\newcommand{\vcf}{\mathrm{VCF}}
\newcommand{\half}{\mathcal{C}^{+}}
\newcommand{\ind}[1]{\mathbf{1}\{#1\}}

\begin{document}

\chapter{Posterior Inference for Treatment Effects}
\label{ch:posterior-inference}

\section{Overview}
\label{sec:pi-overview}

The preceding chapters established a structural decomposition of the observed
outcome into a prognostic baseline, a causal treatment effect, and a
confounding bias term, together with a Bayesian nonparametric prior over the
component functions. That development concerns \emph{how the model is fitted}.
The present chapter concerns \emph{what is done with the fit}: how the posterior
distribution of the component functions is converted into causal conclusions
about treatment effects, how the contribution of the real-world data is
controlled and measured, and how the heterogeneity of both the treatment effect
and the confounding bias is summarised in an interpretable way.

Two estimands are the focus throughout: the conditional average treatment
effect (CATE) and the acceleration factor, its exponential. Both are functionals
of a single component of the decomposition, the treatment-effect function, and
this makes posterior inference for them direct.

The chapter is organised around four tasks. Section~\ref{sec:pi-estimands}
defines the CATE and the acceleration factor from the fitted decomposition and
fixes the target population. Section~\ref{sec:pi-summaries} describes posterior
summaries of conditional and population-averaged effects, including the
treatment of covariate-distribution uncertainty through the Bayesian bootstrap.
Section~\ref{sec:pi-borrowing} introduces the adaptive borrowing architecture:
two scalar parameters that govern, respectively, how strongly the prognostic
surfaces of the two sources are pooled and how strongly the real-world treatment
data are allowed to inform the causal effect.
Section~\ref{sec:pi-diagnostics} develops a suite of diagnostics that report how
much information was borrowed and where. Finally, Section~\ref{sec:pi-vim}
introduces variable importance measures for treatment effect heterogeneity and,
exploiting the decomposition, a novel variable importance measure for the
confounding function itself.

\section{Causal Estimands: the CATE and the Acceleration Factor}
\label{sec:pi-estimands}

\subsection{The decomposition and its parameterisation}

We work with the structural model for the log-survival outcome
$Y = \log T$,
\begin{equation}
\label{eq:decomposition}
Y_i \;=\; m_0(X_i, S_i) \;+\; \tau(X_i)\,A_i \;+\; (1 - S_i)\,A_i\, c(X_i)
        \;+\; \sigma\,\varepsilon_i ,
\end{equation}
where $A \in \{0,1\}$ is the treatment indicator, $S \in \{0,1\}$ is the
data-source indicator with $S = 1$ for the randomized controlled trial (RCT)
and $S = 0$ for the real-world data (RWD), and $X \in \X \subseteq
\mathbb{R}^p$ is a vector of baseline covariates. Here $m_0(\cdot,\cdot)$ is the
baseline (control) outcome surface, $\tau(\cdot)$ is the treatment-effect
function, and $c(\cdot)$ is the confounding function, which enters only for
treated observational units. The residual $\varepsilon$ has a mean-zero,
unit-variance distribution with cumulative distribution function $\Phi$, itself
assigned a nonparametric prior, and $\sigma > 0$ is a scale parameter.

Two modelling choices, used throughout this chapter, refine
\eqref{eq:decomposition}. First, the baseline surface is parameterised as
\begin{equation}
\label{eq:m0-split}
m_0(x, s) \;=\; f(x) \;+\; (1 - s)\, g(x),
\end{equation}
so that $f$ is a baseline surface common to both sources --- interpretable as
the baseline of the RCT population, $s = 1$ --- and $g$ is an RWD-specific
deviation. Second, the confounding function is split into a
covariate-homogeneous intercept and a covariate-varying remainder, in
scale-separated form,
\begin{equation}
\label{eq:c-scale}
c(x) \;=\; c_0 \;+\; \sigma_c\,\tilde c(x),
\end{equation}
where $c_0 \in \mathbb{R}$ is an unshrunk intercept, $\tilde c$ carries a
standard ensemble prior centred at zero, and $\sigma_c \ge 0$ is a global scale.
The intercept $c_0$ absorbs any uniform component of the confounding bias, while
$\sigma_c$ controls the magnitude of its variation across covariate values; the
rationale for this split is given in Section~\ref{subsec:c-interpret}. The
parameters $\sigma_g$ (the leaf scale of $g$) and $\sigma_c$ are the borrowing
controls and are the subject of Section~\ref{sec:pi-borrowing}.

\subsection{The conditional average treatment effect and the acceleration factor}

The primary conditional estimand is the CATE on the log-time scale,
\begin{equation}
\label{eq:cate}
\tau(x) \;=\; \E\bigl[Y(1) - Y(0) \given X = x\bigr]
        \;=\; \E\bigl[\log T(1) - \log T(0) \given X = x\bigr].
\end{equation}
Under the decomposition \eqref{eq:decomposition}, $\tau(x)$ is the
treatment-effect function itself: it is a component of the model rather than a
derived functional, and it does not depend on the baseline surface, the scale
$\sigma$, or the error distribution $\Phi$.

Because the model is of accelerated failure time form, exponentiating the
log-scale effect yields the natural multiplicative estimand. We define the
acceleration factor
\begin{equation}
\label{eq:af}
\af(x) \;=\; \exp\{\tau(x)\}.
\end{equation}
The acceleration factor is the factor by which treatment multiplies survival
time: $\af(x) > 1$ indicates that treatment prolongs survival at covariate value
$x$, and $\af(x) < 1$ that it shortens it. It is reported on the time scale on
which clinical conclusions are usually framed and, like the CATE, is a function
of $\tau$ alone.

\begin{remark}[Robustness to the error model]
\label{rem:phi-free}
Neither \eqref{eq:cate} nor \eqref{eq:af} involves the error distribution
$\Phi$. The nonparametric prior on $\Phi$ governs the shape of the survival
curves and contributes to the efficiency of the fit, but the target estimands
are pure functionals of the treatment-effect function. Conclusions about the
CATE and the acceleration factor are therefore insensitive to the specification
of the error distribution. Should survival-probability contrasts or
restricted-mean survival differences be required as secondary summaries, they
are recovered from the full model through $\Phi$; they are not pursued in this
chapter.
\end{remark}

\subsection{Target population and the role of confounding}
\label{subsec:target-pop}

The causal estimand is the treatment effect in the observational (target)
population. By the cross-source exchangeability assumption the treatment-effect
function $\tau$ does not depend on the data source, so the CATE and the
acceleration factor are common to both populations; what is specific to the
target population is the covariate distribution over which conditional effects
are averaged, addressed in Section~\ref{sec:pi-summaries}.

The confounding function does not enter the causal estimands, but it governs
the relationship between the causal effect and what is \emph{observed} in the
real-world data. In the treated real-world data arm the conditional mean
implied by \eqref{eq:decomposition} is
\[
\E[Y \given X = x, A = 1, S = 0] \;=\; m_0(x, 0) + \tau(x) + c(x),
\]
so the observed treated--control contrast in the real-world data identifies
$\tau(x) + c(x)$, not $\tau(x)$. An analysis of the real-world data alone
would therefore report a confounded CATE $\tau(x) + c(x)$ and, on the
multiplicative scale, a confounded acceleration factor
\[
\exp\{\tau(x) + c(x)\} \;=\; \af(x)\cdot\exp\{c(x)\}.
\]
The confounding thus acts \emph{additively} on the CATE scale and
\emph{multiplicatively} on the acceleration-factor scale, with $\exp\{c(x)\}$
the factor by which the naive observational acceleration factor is inflated. The
role of the decomposition, and of the borrowing architecture of
Section~\ref{sec:pi-borrowing}, is to isolate $\tau(x)$ from $c(x)$ so that the
reported CATE and acceleration factor are the de-confounded, causal quantities.

\subsection{Interpretation of the confounding function}
\label{subsec:c-interpret}

The confounding function is modelled as a function of the observed covariates
$x$, whereas the confounding itself originates in unmeasured covariates $U$ that
are by definition absent from $x$. This invites a natural objection: in giving
$c$ a prior over functions of $x$ alone, is some assumption being smuggled in
about the relationship between $U$ and $X$? The answer is that no such
assumption is being made --- the dependence of $c$ on $x$ is forced by its
definition --- but the structure of $U$ does control how the function behaves,
and this is worth making explicit.

\paragraph{The confounding function is a function of $x$ by construction.}
As defined in the preceding chapter, the confounding function is a contrast of
conditional expectations taken \emph{given} $X = x$; the unmeasured covariates
have already been integrated out in forming those expectations. The result is
therefore necessarily a function of $x$, and modelling it with an ensemble prior
over functions of $x$ is exact rather than approximate. Equivalently, the bias
of the observational treated--control contrast, viewed as a function of $x$, is
\emph{by definition} equal to $c(x)$: the entire confounding bias present in the
conditional mean projects onto a function of the observed covariates, and no
component of it escapes into a direction orthogonal to $x$. The framework does
not assume the confounding to be small, or to be explained by $x$; it represents
it exactly.

\paragraph{The shape of $c$ is governed by interaction between $U$ and $X$.}
What the unmeasured structure controls is not whether $c$ depends on $x$, but
how. Consider the prognostic component of the bias, and write the structural
prognostic function as $f_0(x,u)$ (the prognostic function of the underlying
structural model, distinct from the surface $f$ in \eqref{eq:m0-split}). Let
\[
\Delta(u \given x) \;=\;
  p(u \given x, A{=}1, S{=}0) - p(u \given x, A{=}0, S{=}0)
\]
denote the difference between the treated and control conditional densities of
the unmeasured covariates at covariate value $x$; this difference integrates to
zero over $u$. The prognostic confounding is then
$\int f_0(x,u)\,\Delta(u \given x)\,du$. If the prognostic surface is additive,
$f_0(x,u) = a(x) + b(u)$, the $a(x)$ term integrates against $\Delta$ to zero
and the expression reduces to $\int b(u)\,\Delta(u \given x)\,du$, which depends
on $x$ only through the tilt $\Delta(u \given x)$. If, in addition, that tilt
does not depend on $x$, the confounding is a \emph{constant}. The covariate
dependence of $c(x)$ therefore arises precisely from interaction between $U$ and
$X$ --- either through non-additivity of the prognostic surface in $(x,u)$, or
through covariate dependence of the treatment-assignment mechanism --- and not
from mere marginal association between the two. This is the rationale for the
intercept--remainder split \eqref{eq:c-scale}: the intercept $c_0$ carries the
covariate-homogeneous bias that survives under additive structure, while the
shrunk remainder $\sigma_c \tilde c(x)$ carries the interaction-driven
variation, which is often plausibly small and is therefore an appropriate target
for regularisation.

\paragraph{What the unmeasured covariates still cost.}
Although modelling $c$ as a function of $x$ is exact, three consequences of the
unmeasured confounding remain and should be borne in mind. First, the ensemble
prior and the scale $\sigma_c$ encode a belief that $c$ is regular and ideally
small or smoothly varying; whether that belief holds is itself governed by the
$U$--$X$ interaction structure above, and a confounding mechanism with complex
interaction will be captured poorly, allowing residual bias to leak into $\tau$.
Second, the confounding function conflates two sources of bias: an imbalance in
the baseline prognostic surface across treatment arms, and --- when the
treatment effect itself depends on $U$ --- an imbalance in effect modification.
These two are not separately identified, so $c$ is properly read as the total
observable footprint of confounding rather than as prognostic imbalance alone.
Third, and most consequentially, the separation of $\tau$ from $c$ rests on the
cross-source transportability assumption, which is itself a condition on $U$.

\begin{remark}[Transportability is the binding assumption on $U$]
\label{rem:transport-U}
Identification of $\tau$ and $c$ requires that the conditional law of the
potential outcomes given $X$ be common to the two sources, which in turn
requires the conditional distribution of $U$ given $X$ to be comparable across
the trial and the real-world data. If the two populations differ on an
unmeasured axis, so that $U \given X$ is not transportable, then $\tau(x)$ is no
longer common to the sources and the confounding function will absorb both
genuine within-source confounding and cross-source non-transportability, with no
internal means of distinguishing them. This is the assumption the framework
cannot diagnose from the observed data, and it is where the relationship of the
unmeasured covariates to the data \emph{sources}, rather than to $X$, becomes
critical. Sensitivity analysis with respect to this assumption is an essential
companion to the analysis.
\end{remark}

\section{Posterior Summaries of Treatment Effects}
\label{sec:pi-summaries}

Posterior inference proceeds from $R$ draws of the Markov chain Monte Carlo
sampler. Each draw $r = 1, \dots, R$ supplies a complete set of component
functions and scale parameters,
\[
\bigl\{\, f^{(r)},\; g^{(r)},\; \tau^{(r)},\; c^{(r)},\; \sigma^{(r)},\;
\Phi^{(r)} \,\bigr\},
\]
with $m_0^{(r)}(x, s) = f^{(r)}(x) + (1-s) g^{(r)}(x)$. For the estimands of this
chapter only the draws of $\tau$ are needed.

\subsection{Conditional estimands}

For a fixed covariate value $x$, the posterior of the CATE is read directly off
the treatment-effect draws, $\{\tau^{(r)}(x)\}_{r=1}^{R}$, and the posterior of
the acceleration factor is obtained by the monotone transform
$\af^{(r)}(x) = \exp\{\tau^{(r)}(x)\}$. The posterior mean and equal-tailed
credible intervals follow from the empirical summaries of these draws; because
the transform \eqref{eq:af} is monotone, a credible interval for the
acceleration factor is the exponential of the credible interval for the CATE.

\subsection{Population-averaged estimands}

The population estimand is the average of the conditional effect over the
covariate distribution of the observational population. The population CATE is
\begin{equation}
\label{eq:pop-cate}
\bar\tau \;=\; \E\bigl[\tau(X) \given S = 0\bigr].
\end{equation}
The corresponding population acceleration factor requires a choice, because
exponentiation does not commute with averaging. We define the population
acceleration factor as the exponential of the population CATE,
\begin{equation}
\label{eq:pop-af}
\overline{\af} \;=\; \exp\{\bar\tau\}
              \;=\; \exp\bigl\{\E[\tau(X)\given S=0]\bigr\},
\end{equation}
the geometric mean of the conditional acceleration factors. This is the coherent
multiplicative counterpart of the population CATE and is the summary reported
below. The arithmetic-mean alternative $\E[\exp\{\tau(X)\}\given S=0]$ is also
available from the same draws; it exceeds \eqref{eq:pop-af} by a Jensen gap and
is the appropriate summary only when an average over individuals of the
multiplicative effect itself is the quantity of interest.

\subsection{Covariate-distribution uncertainty}

A plug-in estimate of \eqref{eq:pop-cate} averages $\tau^{(r)}$ over the observed
covariate values $\{x_i : S_i = 0\}$, treating them as fixed. For a model
\emph{parameter} this is appropriate, but $\bar\tau$ is a \emph{functional of the
covariate distribution}, and the observational sample is itself only a finite
draw from the target population; conditioning on the empirical covariate
distribution understates the posterior uncertainty of $\bar\tau$ and of
$\overline{\af}$.

We propagate this additional uncertainty with the Bayesian bootstrap
\citep{rubin1981}, following its use for causal functionals by
\citet{antonelli2019}. A noninformative Dirichlet posterior is placed on the
target covariate distribution: at draw $r$, weights
\[
w^{(r)} = \bigl(w_1^{(r)}, \dots, w_{n_0}^{(r)}\bigr)
\;\sim\; \mathrm{Dirichlet}(1, \dots, 1)
\]
are drawn over the $n_0$ observational units, and the population CATE is formed
as the reweighted average
\begin{equation}
\label{eq:bb-cate}
\bar\tau^{(r)} \;=\; \sum_{i : S_i = 0} w_i^{(r)}\, \tau^{(r)}(x_i),
\qquad
\overline{\af}^{(r)} \;=\; \exp\{\bar\tau^{(r)}\}.
\end{equation}
The resulting posteriors reflect both model uncertainty, through the
treatment-effect draws, and covariate-sampling uncertainty, through the
Dirichlet weights. The procedure is summarised in
Algorithm~\ref{alg:pop-effect}.

\begin{quote}
\noindent\textbf{Algorithm~\ref{alg:pop-effect}: Posterior inference for the
population CATE and acceleration factor.}\refstepcounter{equation}\label{alg:pop-effect}
\begin{enumerate}[label=\textup{(\arabic*)},leftmargin=2.4em]
\item For each posterior draw $r = 1, \dots, R$:
  \begin{enumerate}[label=\textup{(\alph*)},leftmargin=2.0em]
  \item evaluate the treatment-effect function $\tau^{(r)}(x_i)$ at every
        observational unit $i$ with $S_i = 0$;
  \item draw Dirichlet weights $w^{(r)} \sim \mathrm{Dirichlet}(1,\dots,1)$
        over the $n_0$ observational units;
  \item form the reweighted population CATE $\bar\tau^{(r)}$ and the population
        acceleration factor $\overline{\af}^{(r)}$ via \eqref{eq:bb-cate}.
  \end{enumerate}
\item Report the posterior means $R^{-1}\sum_r \bar\tau^{(r)}$ and
      $R^{-1}\sum_r \overline{\af}^{(r)}$, and the empirical quantiles of
      $\{\bar\tau^{(r)}\}_r$ and $\{\overline{\af}^{(r)}\}_r$ as credible
      intervals.
\end{enumerate}
\end{quote}

\section{Adaptive Borrowing of Real-World Data}
\label{sec:pi-borrowing}

The motivation for combining the two data sources is efficiency: although the
CATE is identified from the RCT alone, the real-world data carries
additional information that can sharpen its estimation. The amount of
information transferred --- the \emph{borrowing} --- should not be fixed in
advance but determined by the data. This section describes how borrowing is
parameterised and shows that a single scalar governs how strongly the
observational treatment data inform the causal effect.

\subsection{Two channels of information sharing}

The real-world data contributes to the fit through two distinct channels,
each governed by one of the functions introduced in \eqref{eq:m0-split} and
\eqref{eq:c-scale}.

\paragraph{Prognostic channel.}
Through \eqref{eq:m0-split}, the observational controls inform the shared
baseline $f$, while the source-specific deviation $g$ absorbs any systematic
difference between the two baseline surfaces. Under cross-source
exchangeability, $g$ captures only covariate-shift and measurement-related
discrepancies in the prognostic surface and is expected to be small and smooth.
The leaf-scale parameter $\sigma_g$ controls its magnitude; assigning it a
half-Cauchy prior, $\sigma_g \sim \half(0, \alpha_g)$, allows the data to
interpolate between $\sigma_g \to 0$, under which the prognostic surfaces of the
two sources are fully pooled, and large $\sigma_g$, under which they are
estimated separately.

\paragraph{Confounding channel.}
Through \eqref{eq:c-scale}, the treated observational units inform the
confounding function $c$. Its variation scale receives a half-Cauchy prior,
$\sigma_c \sim \half(0, \alpha_c)$, with $\sigma_c = 0$ corresponding to
confounding that is uniform across covariates --- carried entirely by the
intercept $c_0$ --- and $c_0 = 0$ additionally to its complete absence.

These two channels are deliberately separated: discrepancy in the baseline
prognostic surface and discrepancy due to confounding are conceptually distinct
phenomena and are assigned independent priors. The prognostic channel is where
the real-world data is most useful and least dangerous; the confounding
channel is where its contribution must be regulated.

\subsection{Borrowing on the treatment effect is governed by the confounding scale}

The central structural fact of the decomposition is that the borrowing of
information about the \emph{treatment effect} is controlled entirely by the
shrinkage of the confounding function.

\begin{proposition}[Confounding scale as the borrowing control]
\label{prop:borrowing}
Consider the conditional mean of the treated real-world data arm implied by
\eqref{eq:decomposition},
\[
\E[Y \given X = x, A = 1, S = 0] \;=\; m_0(x, 0) + \tau(x) + c(x).
\]
If $c$ is left unconstrained, the treated observational units identify only the
sum $\tau(x) + c(x)$ and contribute no information about $\tau$ separately; the
CATE is then identified by the RCT alone. As the prior on $c$ shrinks it
towards the zero function, the treated observational units increasingly behave
as unconfounded and contribute directly to the estimation of $\tau$.
Consequently the scale $\sigma_c$ in \eqref{eq:c-scale} acts as the effective
borrowing-strength parameter for the treatment effect: the limit
$\sigma_c \to 0$ corresponds to full pooling of the two sources for $\tau$,
and the limit $\sigma_c \to \infty$ corresponds to RCT-only inference.
\end{proposition}

\noindent
The justification is an identification argument. Within the real-world data
the treated--control contrast identifies $\tau(x) + c(x)$ but not its two
components separately; separation requires the external information supplied by
the RCT. A flexible, unconstrained $c$ therefore absorbs whatever the treated
real-world data would otherwise contribute to $\tau$, leaving the RCT as the
sole source of information about the treatment effect. Shrinking $c$ removes
this absorptive capacity in proportion to $\sigma_c$, so that the observational
treated arm is progressively pooled with the trial.

Proposition~\ref{prop:borrowing} has a useful practical reading. Two analyses a
practitioner might perform without a formal model --- pooling the two datasets
and ignoring confounding, or analysing the trial in isolation --- are precisely
the two endpoints $\sigma_c = 0$ and $\sigma_c \to \infty$. The proposed model
does not commit to either: it places a prior over the continuum between them and
lets the data select the operating point. It also shows why a likelihood-level
downweighting of the real-world data, such as a power prior, is not the
appropriate instrument here: such a device would simultaneously discount the
prognostic channel, where the real-world data are most valuable, whereas the
relevant control acts only on the confounding channel.

\begin{remark}[Effect of the intercept on the borrowing limits]
\label{rem:intercept-borrowing}
Under the intercept--remainder parameterisation \eqref{eq:c-scale}, the limit
$\sigma_c \to 0$ leaves the confounding at the constant $c_0$ rather than at
zero. Because $c_0$ is a single scalar, it is identified from the contrast
between the RCT and the observational treated arm and is estimated regardless of
$\sigma_c$. In this limit the observational treated units therefore still
contribute fully to the \emph{covariate-varying structure} of $\tau$ --- its
heterogeneity --- while a single global offset is removed by $c_0$. The
statement of Proposition~\ref{prop:borrowing} thus refines to: $\sigma_c$ governs
the borrowing of treatment-effect \emph{heterogeneity}, with the
covariate-homogeneous level handled separately by $c_0$. The limit
$\sigma_c \to \infty$ is unchanged, and continues to correspond to RCT-only
inference for $\tau$.
\end{remark}

\subsection{Specification of the borrowing priors}

The half-Cauchy scale $\alpha_c$ encodes the analyst's prior stance on
confounding. A small $\alpha_c$ expresses skepticism that the observational
study is confounded and therefore encourages borrowing; this is the stance of
the commensurate-prior tradition for historical-data borrowing
\citep{hobbs2011}. A larger $\alpha_c$ is agnostic and lets the data determine
the degree of borrowing, in the spirit of data-fusion estimators. Because the
two stances can yield materially different conclusions, both should be
implemented and contrasted, and $\alpha_c$ may be calibrated against the range
of confounding magnitudes judged plausible --- which, on the acceleration-factor
scale, can be elicited as a plausible range for the multiplicative bias
$\exp\{c(x)\}$. The prognostic scale $\alpha_g$ is less delicate, since
exchangeability makes a large prognostic discrepancy implausible, and a weakly
informative choice is generally adequate.

\begin{remark}[Finer, region-specific borrowing]
\label{rem:spike-slab}
A single global scale $\sigma_c$ permits only one borrowing strength for the
whole covariate space. A finer architecture assigns each tree $t$ in the
ensemble for $\tilde c$ an inclusion indicator $\gamma_t \sim
\mathrm{Bernoulli}(\pi)$, with $\gamma_t = 0$ deactivating the tree and
$\pi \sim \mathrm{Beta}(a_\pi, b_\pi)$. The posterior of $\pi$ then reports the
fraction of trees required to represent the confounding, and the posterior of
$\sum_t \gamma_t$ summarises the complexity of the confounding structure. This
extension nests the global-scale specification as the special case
$\gamma_t \equiv 1$ and yields additional diagnostics, at the cost of one
Bernoulli update per tree.
\end{remark}

\section{Diagnostics for Borrowing}
\label{sec:pi-diagnostics}

A combined analysis of randomized and real-world data is only credible if it
can report \emph{how much} information was borrowed, \emph{where} in the
covariate space the borrowing occurred, and \emph{how sensitive} the conclusions
are to that borrowing. Because the borrowing controls of
Section~\ref{sec:pi-borrowing} are explicit model parameters, most diagnostics
reduce to interpretable summaries of those parameters and of the functions they
scale. We organise the suite by the question each diagnostic answers.

\subsection{How much was borrowed: global scalar diagnostics}

\paragraph{Posterior of the borrowing scales.}
The posteriors of $\sigma_c$ and $\sigma_g$ are the most direct summaries. To be
interpretable they should be reported in meaningful units. For the confounding
channel, summarising the multiplicative bias $\exp\{c(X)\}$ --- for example its
posterior mean and a central interval over the target population --- expresses
the confounding on the same acceleration-factor scale as the estimand, in the
form ``the confounding inflates the observational acceleration factor by a
factor between $b_{\text{low}}$ and $b_{\text{high}}$''. The split
\eqref{eq:c-scale} additionally allows the homogeneous and heterogeneous parts
of the confounding to be reported separately: the intercept $c_0$ as a single
uniform bias, and the posterior spread of $\sigma_c \tilde c(X)$ as the extent
to which that bias varies across patients. The scale $\sigma_g$ is summarised
analogously as a prognostic discrepancy between the two sources.

\paragraph{Effective sample size for the treatment effect.}
The practical value of borrowing is the variance reduction it produces. Let
$\Var_{\text{full}}\{\tau(x) \given \data\}$ denote the posterior variance of
the CATE under the combined model and $\Var_{\text{rct}}\{\tau(x) \given \data\}$
the corresponding variance under a refit using the RCT alone. The borrowing is
then expressed as an effective number of additional trial patients,
\begin{equation}
\label{eq:ess}
\mathrm{ESS}_{\text{add}}(x)
  \;=\; n_{\text{rct}}
  \left\{
  \frac{\Var_{\text{rct}}\{\tau(x) \given \data\}}
       {\Var_{\text{full}}\{\tau(x) \given \data\}} - 1
  \right\}.
\end{equation}
Averaged over the target covariate distribution, \eqref{eq:ess} provides a
single headline figure --- the borrowing was equivalent to enrolling
$\mathrm{ESS}_{\text{add}}$ additional randomized patients --- in the idiom
established for historical-data borrowing
\citep{neuenschwander2010,hupf2021}. Its evaluation requires one RCT-only refit
alongside the main analysis. The effective sample size is computed on the CATE
scale, on which the variance is naturally defined; the acceleration factor,
being a monotone transform, inherits the same conclusion.

\subsection{Where it was borrowed: functional diagnostics}

Scalar summaries average over the covariate space and conceal heterogeneity in
the borrowing. Two functional diagnostics localise it.

\paragraph{Borrowing map.}
The pointwise posterior of the confounding function --- for instance the
posterior mean $\E\bigl[|c(x)| \given \data\bigr]$, or the exceedance
probability $\Prob\bigl[|c(x)| > \delta \given \data\bigr]$ for a clinically
meaningful threshold $\delta$ --- displayed over the covariate space, identifies
the regions in which the model infers confounding. When $p$ is large the map is
shown over the leading principal components or conditional on a small set of
clinically salient covariates. Presenting the posterior mean alongside the
credible-interval width separates regions of inferred confounding from regions
of mere uncertainty.

\paragraph{Local effective sample size.}
Evaluating \eqref{eq:ess} as a function of $x$ produces a map of where the
real-world data contributed and where it did not. Plotting
$\mathrm{ESS}_{\text{add}}(x_i)$ against the inferred local bias $|c(x_i)|$
exhibits the adaptive behaviour the framework is designed to deliver: regions of
small inferred confounding receive substantial borrowing, regions of large
inferred confounding receive little, and the scatter is itself a visual proof of
that mechanism.

\subsection{Whether to trust it: robustness checks}

\paragraph{Prior--data conflict.}
Overlaying the prior and posterior densities of $\sigma_c$, and of $\sigma_g$,
indicates whether the data were informative about the degree of confounding.
Near-coincident densities signal that the data carry little evidence either way;
a posterior concentrated away from the prior indicates that the data have
spoken. Formal prior--data conflict checks \citep{evans2006} can supplement the
graphical comparison.

\paragraph{Tipping-point analysis.}
Refitting the model with the borrowing scale fixed at a grid of values, from
$\sigma_c = 0$ (full pooling) through the data-driven posterior median to a large
value approximating RCT-only inference, and tracing the population CATE or
acceleration factor along that grid, reveals the dependence of the conclusion on
the degree of borrowing. The data-driven fit should lie between the extremes and
the estimand should vary smoothly; an abrupt change, or a reversal of the
substantive conclusion in the neighbourhood of the operating point, identifies a
fragile result.

\paragraph{Cross-source posterior predictive checks.}
Simulating treated observational outcomes from the posterior predictive
distribution and comparing them with the observed treated observational
outcomes --- for example through an overlay of survival curves with simulation
envelopes \citep{gelman1996} --- provides a model-free check on the confounding
correction. Systematic discrepancy confined to the treated arm is the signature
of confounding that the model has failed to absorb.

Table~\ref{tab:diag} collects the suite.

\begin{table}[t]
\centering
\caption{Borrowing diagnostics organised by the question addressed.}
\label{tab:diag}
\begin{tabular}{@{}lll@{}}
\toprule
Question & Diagnostic & Form \\
\midrule
How much was borrowed?
  & Posterior of $\sigma_c, \sigma_g$ & global scalar \\
  & Effective sample size, Eq.~\eqref{eq:ess} & global scalar \\[2pt]
Where was it borrowed?
  & Borrowing map & function of $x$ \\
  & Local effective sample size & function of $x$ \\[2pt]
Can it be trusted?
  & Prior--data conflict & prior vs.\ posterior \\
  & Tipping-point analysis & estimand vs.\ $\sigma_c$ \\
  & Cross-source predictive check & survival-curve overlay \\
\bottomrule
\end{tabular}
\end{table}

\section{Variable Importance for Heterogeneity and Confounding}
\label{sec:pi-vim}

A flexible estimate of the CATE describes treatment effect heterogeneity but
does not, on its own, identify which covariates drive it. Treatment effect
variable importance measures \citep{hines2022} summarise the contribution of a
subset of covariates to the heterogeneity of the CATE through a model-free,
nonparametric quantity with an analysis-of-variance interpretation. This section
adapts that idea to posterior inference within the present framework and, using
the decomposition, extends it to the confounding function.

\subsection{The variance of the treatment effect and the TE-VIM}

All expectations and variances in this section are taken over the covariate
distribution of the target (observational) population. The measures are computed
on the CATE scale, $\tau$, on which the model is additive and the
variance decomposition carries its clean analysis-of-variance interpretation;
the same covariates govern the heterogeneity of the acceleration factor, which
is a monotone transform of $\tau$. The global measure of heterogeneity is the
variance of the treatment effect,
\begin{equation}
\label{eq:vte}
\vte \;=\; \Var\{\tau(X)\},
\end{equation}
which quantifies how much the treatment effect varies across the population
\citep{levy2021}.

\begin{definition}[Treatment effect variable importance measure]
\label{def:tevim}
For a subset $s \subseteq \{1, \dots, p\}$ of covariate indices, let $X_{-s}$
denote the covariates outside $s$ and define the reduced CATE
\[
\tau_{-s}(x_{-s}) \;=\; \E\bigl[\tau(X) \given X_{-s} = x_{-s}\bigr].
\]
The treatment effect variable importance measure of $s$ is
\begin{equation}
\label{eq:tevim}
\theta_s \;=\; \E\bigl[\{\tau(X) - \tau_{-s}(X_{-s})\}^2\bigr]
        \;=\; \vte - \Var\{\tau_{-s}(X_{-s})\},
\end{equation}
and its standardised form is the proportion $\psi_s = \theta_s / \vte$.
\end{definition}

\noindent
The two expressions in \eqref{eq:tevim} coincide because $\tau_{-s}$ is the
$L^2$ projection of $\tau$ onto functions of $X_{-s}$. The measure $\theta_s$ is
the increase in mean-squared error incurred when the covariates in $s$ are
removed from the conditioning set, equivalently the portion of treatment effect
heterogeneity attributable to $s$; the standardised $\psi_s$ reads as the share
of $\vte$ explained by $s$, giving the analysis-of-variance interpretation. The
construction has been extended to censored survival outcomes by
\citet{ziersen2024}, with which the present log-scale CATE is consistent.

\subsection{Posterior inference for variable importance}
\label{subsec:vim-posterior}

The variable importance measure of Definition~\ref{def:tevim} is a functional of
the CATE, and within the present framework its posterior is obtained directly,
without the influence-function machinery and sample splitting required by
semiparametric estimators. For each posterior draw $r$:
\begin{enumerate}[label=\textup{(\arabic*)},leftmargin=2.4em]
\item evaluate the treatment-effect function $\tau^{(r)}(x_i)$ at every target
      unit;
\item estimate the reduced CATE $\tau^{(r)}_{-s}(x_{-s}) =
      \E\bigl[\tau^{(r)}(X) \given X_{-s} = x_{-s}\bigr]$, the conditional
      expectation of the draw's treatment-effect function given the retained
      covariates, by a flexible nonparametric regression of
      $\{\tau^{(r)}(x_i)\}_i$ on $\{x_{-s,i}\}_i$ --- equivalently, by
      marginalising $\tau^{(r)}$ over the conditional covariate distribution of
      $X_s$ given $X_{-s}$;
\item form
      $\theta_s^{(r)} = \sum_i w_i^{(r)}
      \bigl\{\tau^{(r)}(x_i) - \tau^{(r)}_{-s}(x_{-s,i})\bigr\}^2$,
      with $w^{(r)}$ the Dirichlet weights of \eqref{eq:bb-cate}.
\end{enumerate}
The collection $\{\theta_s^{(r)}\}_r$ is the posterior of the variable
importance measure, from which the posterior of $\psi_s$ and its credible
interval follow. The Dirichlet weights again propagate covariate-distribution
uncertainty, so the posterior reflects both model and sampling variability. This
posterior-based route is also a substantive advantage: influence-function
estimators of variable importance are known to have unreliable calibration at
realistic sample sizes, whereas the posterior delivers uncertainty
quantification intrinsically.

\begin{remark}[The reduction is a conditional mean, not a restricted projection]
\label{rem:cond-mean}
The reduced function in Definition~\ref{def:tevim} is the conditional
expectation $\E[\tau(X) \given X_{-s}]$, equivalently the $L^2$-optimal
predictor of $\tau(X)$ among \emph{all} functions of $X_{-s}$. Step~(2) must
therefore estimate it with a flexible learner. Restricting the reduction to a
linear --- or otherwise low-complexity --- class would change the estimand from
``the heterogeneity explained by $X_{-s}$'' into ``the heterogeneity explained
by a function of $X_{-s}$ within that restricted class'', understating
importance and potentially reordering the covariate ranking, the more so because
the treatment-effect function is itself deliberately nonlinear. The
conditional-mean reduction is the convention of \citet{hines2022} and of the
multivariate-treatment importance measures of \citet{shin2024}. A deliberate
projection of the CATE onto a linear or sparse summary is a legitimate but
\emph{distinct} object --- the posterior projection of \citet{woody2020}, or the
best partially linear projection of \citet{ziersen2024} --- which answers a
different question and is not a substitute for the conditional-mean reduction
used here.
\end{remark}

\subsection{Strategies for estimating the variable importance measures}
\label{subsec:vim-estimation}

The variable importance measures of Definitions~\ref{def:tevim}
and~\ref{def:cvim} are population functionals, and several estimation strategies
for them appear in the literature. They differ along three axes: how the CATE is
estimated, how the reduced conditional mean is obtained, and how the importance
functional is assembled and debiased. We summarise the main options and situate
the posterior plug-in of Section~\ref{subsec:vim-posterior} among them; the same
options apply to the confounding measure with $c$ in place of $\tau$.

\paragraph{One-step efficient-influence-curve estimator.}
The semiparametric route of \citet{hines2022} constructs the augmented
inverse-probability-weighted pseudo-outcome $\varphi(Z)$, which behaves in
expectation like the unobserved individual treatment effect, and estimates
\[
\widehat\theta_s = n^{-1}\sum_{i=1}^n
\Bigl[ \{\widehat\varphi(z_i) - \widehat\tau_s(x_i)\}^2
      - \{\widehat\varphi(z_i) - \widehat\tau(x_i)\}^2 \Bigr].
\]
This is a one-step (estimating-equation) estimator built from the efficient
influence curve of $\theta_s$: it is Neyman-orthogonal, supplies a closed-form
influence-curve variance, and is paired with cross-fitting to avoid
empirical-process (Donsker) conditions. It is not a plug-in estimator and may
return negative values. The CATE $\widehat\tau$ may be supplied by a T-learner
or, for robustness, by the doubly robust DR-learner, and the reduced CATE
$\widehat\tau_s$ by regressing either $\widehat\tau$ or $\widehat\varphi$ on
$X_{-s}$.

\paragraph{Plug-in and targeted (TMLE) estimators.}
A direct plug-in $\widehat\theta_s = \widehat{\Var}\{\widehat\tau(X)\} -
\widehat{\Var}\{\widehat\tau_s(X)\}$ respects the parameter space but inherits
the regularisation bias of the nuisance estimators. The debiased plug-in is the
targeted maximum likelihood estimator, in which the nuisance fits are updated so
that the plug-in is also efficient: \citet{levy2021} give such an estimator for
the VTE and \citet{li2023} for the TE-VIM. Targeting $\theta_s$ is more delicate
than targeting the VTE, because the update must move the full and reduced
regressions compatibly.

\paragraph{Bayesian posterior plug-in.}
When a full posterior over the outcome and treatment-effect functions is
available, as in the present framework, the importance measure is a
deterministic functional of each posterior draw, and its posterior is obtained
by evaluating that functional draw by draw --- the route of \citet{shin2024} and
the one adopted in Section~\ref{subsec:vim-posterior}. Three features make it the
natural choice here. First, uncertainty quantification comes directly from the
posterior, with no influence curve and no cross-fitting. Second, the reduced and
full CATE are automatically \emph{compatible} within a draw: because
$\tau^{(r)}_{-s}$ is the conditional mean of the same drawn function
$\tau^{(r)}$, the identity $\E\{\tau\} = \E\{\tau_{-s}\}$ holds exactly at every
draw --- a property that \citet{hines2022} must engineer through their choice of
reduced regression. Third, the regularisation is supplied by the model's own
priors; in particular, the shrinkage on the heterogeneity components
(Section~\ref{sec:pi-borrowing}) pulls the importance measures toward zero when
heterogeneity or confounding is absent. The validity of the posterior plug-in
rests on the outcome model being well-calibrated; \citet{shin2024} establish
that posterior contraction for the outcome surface transfers to the importance
measures.

\paragraph{Separate-sample estimation for inference at the null.}
At the zero-importance boundary $\theta_s = 0$ the influence curve of the
one-step estimator degenerates and asymptotic normality fails
(Remark~\ref{rem:null}). \citet{hines2022} restore valid inference by estimating
$\Var\{\tau(X)\}$ and $\Var\{\tau_s(X)\}$ with efficient estimators on
\emph{independent} subsamples and differencing them, at the cost of efficiency.
The posterior analogue is to avoid testing $\theta_s = 0$ directly and instead
compare importances through the posterior of a \emph{difference},
$\psi_{\{j\}} - \psi_{\{k\}}$, which is well behaved away from the joint
boundary; this is the testing device of \citet{shin2024} and the one we use for
the effect-modifier-versus-confounder comparisons of
Section~\ref{subsec:effmod}.

\paragraph{Computational devices.}
Two accelerations are worth noting. Cross-fitting (sample splitting) is required
for the one-step estimator and optional but stabilising for the others. For the
posterior plug-in, the bottleneck is the per-draw evaluation of the reduced
conditional mean over all units, which amounts to an $n \times n$ evaluation of
the treatment-effect function at each draw. The blocking scheme of
\citet{shin2024} partitions the sample into $K$ blocks, computes the importance
measure within each block, and averages across blocks; it is near-identical to
the full computation in their simulations and is substantially faster. For a
data-fusion analysis with a large observational registry, the blocking scheme is
the practical default.

\begin{remark}[Choosing the covariate subsets]
\label{rem:vim-modes}
Distinct from how $\theta_s$ is estimated is the choice of which subsets $s$ to
evaluate. \citet{hines2022} identify four modes: leave-one-out ($s$ a single
covariate), keep-one-in ($s$ all but one covariate), Shapley aggregation over all
subsets, and covariate grouping guided by domain knowledge. When covariates are
correlated, single-covariate (leave-one-out) importances are diluted --- two
correlated effect modifiers each receive a small measure --- so keep-one-in or
covariate grouping is preferable in that case.
\end{remark}

\subsection{A variable importance measure for confounding}

The decomposition isolates a second heterogeneity function, the confounding
function $c$, and the construction of Definition~\ref{def:tevim} applies to it
verbatim. Let $\vcf = \Var\{c(X)\}$ be the total variation of the confounding
bias across the target population, and for a subset $s$ define the reduced
confounding function $c_{-s}(x_{-s}) = \E[c(X) \given X_{-s} = x_{-s}]$. Under
the parameterisation \eqref{eq:c-scale} the homogeneous intercept $c_0$ cancels
from every variance and centred contrast below, so the confounding variable
importance measure speaks only to the covariate-varying part of the
confounding --- which is appropriate, since a uniform bias is by construction
not attributable to any covariate.

\begin{definition}[Confounding variable importance measure]
\label{def:cvim}
The confounding variable importance measure of a covariate subset $s$ is
\begin{equation}
\label{eq:cvim}
\theta^{c}_s \;=\; \E\bigl[\{c(X) - c_{-s}(X_{-s})\}^2\bigr]
            \;=\; \vcf - \Var\{c_{-s}(X_{-s})\},
\end{equation}
with standardised form $\psi^{c}_s = \theta^{c}_s / \vcf$.
\end{definition}

\noindent
The measure $\theta^{c}_s$ quantifies how much of the confounding bias in the
real-world data is attributable to the covariates in $s$. It answers a
question that has no counterpart in the standard variable importance literature
--- which observed covariates drive the unmeasured confounding --- and is
definable only because the decomposition exposes $c$ as an estimable function.
Its posterior is computed exactly as in the previous subsection, with $c^{(r)}$
in place of $\tau^{(r)}$.

\subsection{Separating effect modifiers from confounding drivers}
\label{subsec:effmod}

Computing the single-covariate importances $\psi_{\{j\}}$ and $\psi^{c}_{\{j\}}$
for each covariate $j$ and plotting them against one another classifies the
covariates into four regimes: pure effect modifiers, important for the treatment
effect but not for the bias; pure confounding drivers, important for the bias
but not for the effect; covariates important for both; and covariates important
for neither. This separation is possible only because the decomposition carries
$\tau$ and $c$ as distinct functions, and it directly serves the applied reader:
it distinguishes covariates that genuinely modify the treatment effect from
covariates that merely reflect how patients were selected into treatment in the
real-world data, a distinction that an analysis of the real-world data
alone cannot draw.

\begin{remark}[The null case]
\label{rem:null}
When a covariate subset has no role in the heterogeneity of interest, the
corresponding importance measure equals zero, which lies on the boundary of the
parameter space. Inference is delicate at this boundary for any estimator. The
measures are therefore best reported as a ranking accompanied by posterior
credible intervals, with values near zero treated as a boundary case rather than
as the basis for a formal test of no effect.
\end{remark}

\subsection{Variable importance as a borrowing diagnostic}

The variable importance measures also feed back into the diagnostic suite of
Section~\ref{sec:pi-diagnostics}. Computing $\psi_{\{j\}}$ for the treatment
effect under the RCT-only refit and under the combined fit shows whether
borrowing sharpens the identification of effect modifiers: tighter posteriors
under the combined fit indicate a genuine gain, while a marked change in the
ranking of effect modifiers is a warning sign, signalling either informative
borrowing or contamination of the trial signal by the real-world data and
warranting closer inspection through the tipping-point analysis.

\section{Summary}
\label{sec:pi-summary}

This chapter has set out how the fitted decomposition is turned into inference
about treatment effects. The two estimands of interest, the conditional average
treatment effect and the acceleration factor, are functionals of the
treatment-effect function alone and are insensitive to the error model; their
posteriors are read directly off the treatment-effect draws and summarised with
the Bayesian bootstrap to account for covariate-distribution uncertainty. The
contribution of the real-world data is governed by two explicit scale
parameters, of which the confounding scale $\sigma_c$ is shown to be the
effective borrowing control for the treatment effect, with the confounding
acting multiplicatively on the acceleration-factor scale. A suite of diagnostics
reports how much was borrowed, where, and how robustly; and variable importance
measures summarise the heterogeneity of both the treatment effect and the
confounding function, the latter being a measure made possible by the
decomposition. Together these elements make the borrowing transparent and
auditable, which is a precondition for the use of real-world data in
confirmatory treatment effect estimation.

% =============================================================================
% Bibliography. When this chapter is merged into the master document, delete
% this block and ensure the cited keys resolve against the master .bib file.
% =============================================================================
\begin{thebibliography}{99}

\bibitem[Antonelli et al.(2019)]{antonelli2019}
Antonelli, J., Parmigiani, G., and Dominici, F. (2019).
High-dimensional confounding adjustment using continuous spike and slab priors.
\textit{Bayesian Analysis}, 14(3), 805--828.

\bibitem[Evans and Moshonov(2006)]{evans2006}
Evans, M. and Moshonov, H. (2006).
Checking for prior-data conflict.
\textit{Bayesian Analysis}, 1(4), 893--914.

\bibitem[Gelman et al.(1996)]{gelman1996}
Gelman, A., Meng, X.-L., and Stern, H. (1996).
Posterior predictive assessment of model fitness via realized discrepancies.
\textit{Statistica Sinica}, 6(4), 733--760.

\bibitem[Hines et al.(2022)]{hines2022}
Hines, O., D\'{i}az-Ordaz, K., and Vansteelandt, S. (2022).
Variable importance measures for heterogeneous treatment effects.
\textit{arXiv preprint} arXiv:2204.06030.

\bibitem[Hobbs et al.(2011)]{hobbs2011}
Hobbs, B. P., Carlin, B. P., Mandrekar, S. J., and Sargent, D. J. (2011).
Hierarchical commensurate and power prior models for adaptive incorporation of
historical information in clinical trials.
\textit{Biometrics}, 67(3), 1047--1056.

\bibitem[Hupf et al.(2021)]{hupf2021}
Hupf, B., Bunn, V., Lin, J., and Dong, C. (2021).
Bayesian semiparametric meta-analytic-predictive prior for historical control
borrowing in clinical trials.
\textit{Statistics in Medicine}, 40(15), 3445--3463.

\bibitem[Levy et al.(2021)]{levy2021}
Levy, J., van der Laan, M., Hubbard, A., and Pirracchio, R. (2021).
A fundamental measure of treatment effect heterogeneity.
\textit{Journal of Causal Inference}, 9(1), 83--108.

\bibitem[Li et al.(2023)]{li2023}
Li, H., Hubbard, A., and van der Laan, M. (2023).
Targeted learning on a variable importance measure for heterogeneous treatment
effects.
\textit{arXiv preprint} arXiv:2309.13324.

\bibitem[Neuenschwander et al.(2010)]{neuenschwander2010}
Neuenschwander, B., Capkun-Niggli, G., Branson, M., and Spiegelhalter, D. J.
(2010).
Summarizing historical information on controls in clinical trials.
\textit{Clinical Trials}, 7(1), 5--18.

\bibitem[Rubin(1981)]{rubin1981}
Rubin, D. B. (1981).
The Bayesian bootstrap.
\textit{The Annals of Statistics}, 9(1), 130--134.

\bibitem[Shin et al.(2024)]{shin2024}
Shin, H., Linero, A., Audirac, M., Irene, K., Braun, D., and Antonelli, J.
(2024).
Treatment effect heterogeneity and importance measures for multivariate
continuous treatments.
\textit{arXiv preprint} arXiv:2404.09126.

\bibitem[Woody et al.(2020)]{woody2020}
Woody, S., Carvalho, C. M., and Murray, J. S. (2020).
Bayesian inference with posterior projection.
\textit{Journal of Computational and Graphical Statistics}, 29(4), 798--808.

\bibitem[Ziersen et al.(2024)]{ziersen2024}
Ziersen, S. C. et al. (2024).
Variable importance measures for heterogeneous treatment effects with survival
outcome.
\textit{arXiv preprint} arXiv:2412.11790.

\end{thebibliography}

\end{document}